\usepackage{hyperref}
\renewcommand{\thesubsubsection}{\Alph{subsubsection}}
\section*{Suggested solutions}
These solutions were prepared by \href{https://openai.com/index/gpt-6-astra/}{ChatGPT Astra 6 from OpenAI} for this TP Web Lab adaptation. 
They are not an official answer key by Pascal Michaillat. 
They accompany the original \href{ps1.tex}{Problem Set on Dynamic Programming}. 
The question letters match the exercise page.

We impose positive consumption, nonnegative capital, feasible asset positions and the relevant transversality conditions. First-order equalities below concern interior optima with differentiable value functions; boundary choices require inequalities. In Problem 2 we read the source's incomplete probability bound as $0.5<\rho<1$ and assume $0<\beta<1$.

\subsection*{Problem 1: optimal growth}
\setcounter{subsubsection}{0}
\subsubsection{Lagrangian and consumption dynamics}
Use current-value multipliers $\lambda_t$ and write
\[
\mathcal{L}=\sum_{t=0}^{\infty}\beta^t\left[\ln c_t+\lambda_t(a k_t^{\alpha}-c_t-k_{t+1})\right].
\]
The consumption and capital conditions are
\[
\lambda_t=\frac{1}{c_t},\qquad
\lambda_t=\beta\lambda_{t+1}\alpha a k_{t+1}^{\alpha-1}.
\]
Consequently,
\[
c_{t+1}=\beta\alpha a k_{t+1}^{\alpha-1}c_t.
\]
This Euler equation alone does not determine the initial consumption choice. It must be combined with feasibility and a transversality condition, here
\[
\lim_{T\to\infty}\beta^T\frac{k_{T+1}}{c_T}=0.
\]
Part F supplies the optimal policy and thus the complete consumption path.

\subsubsection{State and control}
Capital $k_t$ is the state. Consumption $c_t$ is a control; equivalently, choose next-period capital $k_{t+1}$ and recover consumption from the resource constraint.

\subsubsection{Bellman equation}
\[
V(k)=\max_{0\leq k'<a k^{\alpha}}\left[\ln(a k^{\alpha}-k')+\beta V(k')\right].
\]
For the infinite-horizon problem, choosing $k'=0$ makes positive consumption impossible in the following period. The optimum therefore has $k'>0$.

\subsubsection{Euler equation from the Bellman equation}
The first-order and envelope conditions give
\[
\frac{1}{c}=\beta V'(k'),\qquad V'(k)=\frac{\alpha a k^{\alpha-1}}{c}.
\]
Applying the envelope condition next period and substituting into the first-order condition yields exactly the Euler equation in part A.

\subsubsection{First two value-function iterates}
Starting from $V_0=0$, the one-period problem consumes all output and leaves zero terminal capital:
\[
V_1(k)=\ln(a k^{\alpha})=\ln a+\alpha\ln k.
\]
The second iterate solves
\[
V_2(k)=\max_{0<k'<a k^{\alpha}}\left[\ln(a k^{\alpha}-k')+\beta\ln a+\alpha\beta\ln k'\right].
\]
Its first-order condition implies
\[
k'_2(k)=\frac{\alpha\beta}{1+\alpha\beta}a k^{\alpha},\qquad
c_2(k)=\frac{a k^{\alpha}}{1+\alpha\beta}.
\]
Therefore
\[
V_2(k)=\ln\left(\frac{a k^{\alpha}}{1+\alpha\beta}\right)+\beta\ln a
+\alpha\beta\ln\left(\frac{\alpha\beta a k^{\alpha}}{1+\alpha\beta}\right).
\]
The zero terminal capital choice for $V_1$ is a finite-horizon boundary condition, not the infinite-horizon policy.

\subsubsection{Infinite-horizon policy}
The proposed value function has derivative
\[
V'(k)=\frac{\alpha}{(1-\alpha\beta)k}.
\]
Substitute it into the first-order condition:
\[
\frac{1}{a k^{\alpha}-k'}=\frac{\alpha\beta}{(1-\alpha\beta)k'}.
\]
Solving gives
\[
k'(k)=\alpha\beta a k^{\alpha},\qquad c(k)=(1-\alpha\beta)a k^{\alpha}.
\]
For completeness, substitution into the Bellman equation determines the constant:
\[
\kappa=\frac{1}{1-\beta}\left[\ln((1-\alpha\beta)a)+\frac{\alpha\beta}{1-\alpha\beta}\ln(\alpha\beta a)\right].
\]
The policy is feasible and positive, satisfies the Euler equation, and satisfies transversality because $\beta^T k_{T+1}/c_T=\beta^T\alpha\beta/(1-\alpha\beta)\to0$.

\subsubsection{Comparing the approaches}
Dynamic programming yields a policy for every state, makes the continuation value explicit and supports numerical iteration and stochastic extensions. Its costs include storing a value function over the state space and establishing the conditions for convergence; the unbounded logarithmic payoff requires care with general convergence claims.

The Lagrangian gives Euler equations and marginal trade-offs directly and can be simpler for analytical or local calculations. However, these conditions must be completed by feasibility, boundary and transversality conditions; they do not by themselves select the optimal path or provide a global policy.

\subsection*{Problem 2: asset pricing}
\setcounter{subsubsection}{0}
\subsubsection{State and control}
The state is $(s,d)$: current asset holdings and the observed dividend. Price is $p(d)$, so it is not an additional independent state. Choose next-period holdings $s'$; consumption is $c=(d+p(d))s-p(d)s'$. Choices are adapted to currently available information.

\subsubsection{Bellman equation}
Let $P_{ll}=P_{hh}=\rho$ and $P_{lh}=P_{hl}=1-\rho$. For $i\in\lbrace l,h\rbrace $,
\[
V(s,d_i)=\max_{s'\in\Gamma(s,d_i)}\left[u((d_i+p_i)s-p_i s')+\beta\sum_{j\in\lbrace l,h\rbrace }P_{ij}V(s',d_j)\right].
\]
Here $p_i=p(d_i)$ and $\Gamma$ denotes feasible positions, including any borrowing or short-sale restrictions and positive consumption. An interior solution is assumed in part C.

\subsubsection{Euler equation}
The first-order condition and envelope condition are
\[
p_i u'(c)=\beta\sum_jP_{ij}V_s(s',d_j),\qquad
V_s(s,d_i)=(d_i+p_i)u'(c).
\]
Combining them gives
\[
p_tu'(c_t)=\beta\mathbb{E}\left[(d_{t+1}+p_{t+1})u'(c_{t+1})\mid s_t,d_t\right].
\]
Along the chosen policy, current holdings and choices are understood as known. With that convention the exercise writes the conditioning simply as $d_t$. At a binding trading constraint the first-order equality need not hold.

\subsubsection{Linear utility and high-dividend prices}
When $u(c)=c$, marginal utility is one. For an unconstrained marginal trade the pricing equations become
\[
\begin{aligned}
p_l&=\beta[\rho(d_l+p_l)+(1-\rho)(d_h+p_h)],\\
p_h&=\beta[(1-\rho)(d_l+p_l)+\rho(d_h+p_h)].
\end{aligned}
\]
Subtracting and setting $\Delta p=p_h-p_l$, $\Delta d=d_h-d_l$ gives
\[
\Delta p=\beta(2\rho-1)(\Delta d+\Delta p),\qquad
\Delta p=\frac{\beta(2\rho-1)}{1-\beta(2\rho-1)}\Delta d>0.
\]
The numerator is positive because $\rho>1/2$ and $d_h>d_l$; the denominator is positive because $\beta<1$ and $2\rho-1<1$.

For an explicit check, define $\bar d=(d_l+d_h)/2$. Adding the two equations gives $\bar p=\beta\bar d/(1-\beta)$, and $p_l=\bar p-\Delta p/2$, $p_h=\bar p+\Delta p/2$. These are the stationary finite prices. The restriction $p=p(d)$ rules out an additional time-varying bubble. Linear utility does not itself select a unique individual consumption plan.

\subsection*{Problem 3: consumption and labor}
\setcounter{subsubsection}{0}
\subsubsection{Deterministic state and controls}
Capital $k$ is the state; next-period capital $k'$ and labor $l\in[0,1]$ are controls. Consumption is $c=f(k,l)-k'$.

\subsubsection{Deterministic Bellman equation}
\[
V(k)=\max_{\substack{k'\geq0,\ 0\leq l\leq1\\ f(k,l)-k'>0}}
\left[u(f(k,l)-k',l)+\beta V(k')\right].
\]

\subsubsection{Optimality conditions}
At an interior choice, differentiating with respect to $k'$ and $l$ and using the envelope condition yields
\[
u_c(c,l)=\beta V'(k'),\qquad
u_c(c,l)f_l(k,l)+u_l(c,l)=0,\qquad
V'(k)=u_c(c,l)f_k(k,l).
\]
Thus
\[
\begin{aligned}
u_c(c_t,l_t)&=\beta u_c(c_{t+1},l_{t+1})f_k(k_{t+1},l_{t+1}),\\
u_c(c_t,l_t)f_l(k_t,l_t)&=-u_l(c_t,l_t).
\end{aligned}
\]
For labor at $l=0$, the marginal expression $u_cf_l+u_l$ must be nonpositive; at $l=1$ it must be nonnegative. These are one-sided necessary conditions.

The statement assumes utility is convex in labor. This does not ensure a concave maximization problem or an interior labor optimum. Accordingly, the equalities are necessary conditions for an interior optimum, not a general sufficiency theorem. Candidate solutions must be compared with boundary choices.

\subsubsection{Steady-state ratios with Cobb--Douglas production}
Assume a positive stationary solution, with $0<\alpha<1$, $k'=k$ and constant consumption and labor. The capital Euler equation gives
\[
1=\beta\alpha k^{\alpha-1}l^{1-\alpha}
=\beta\alpha(l/k)^{1-\alpha}.
\]
Hence
\[
\frac{l}{k}=\left(\frac{1}{\alpha\beta}\right)^{1/(1-\alpha)},\qquad
\frac{c}{k}=\frac{f(k,l)}{k}-1=\frac{1}{\alpha\beta}-1.
\]
The subtraction of one follows from the exercise's law of motion, which has no undepreciated capital term. These ratios do not determine the levels: the utility function and the labor condition determine the scale, subject to $l\leq1$. Their derivation is conditional on the existence of a positive steady state.

\subsubsection{Stochastic Bellman equation}
The state is now $(k,a_i)$, with the productivity shock observed before choices. Using $P_{ii}=\rho$ and $P_{ij}=1-\rho$ for $j\ne i$,
\[
V(k,a_i)=\max_{\substack{k'\geq0,\ 0\leq l\leq1\\ a_i f(k,l)-k'>0}}
\left[u(a_i f(k,l)-k',l)+\beta\sum_{j=1}^2P_{ij}V(k',a_j)\right].
\]

\subsubsection{Stochastic optimality conditions}
The first-order and envelope conditions at interior choices are
\[
\begin{aligned}
u_c(c,l)&=\beta\sum_jP_{ij}V_k(k',a_j),\\
u_c(c,l)a_i f_l(k,l)+u_l(c,l)&=0,\\
V_k(k,a_i)&=u_c(c,l)a_i f_k(k,l).
\end{aligned}
\]
Combining them yields
\[
\begin{aligned}
u_c(c_t,l_t)&=\beta\mathbb{E}_t\left[u_c(c_{t+1},l_{t+1})a_{t+1}f_k(k_{t+1},l_{t+1})\right],\\
u_c(c_t,l_t)a_t f_l(k_t,l_t)&=-u_l(c_t,l_t).
\end{aligned}
\]
Future consumption and labor depend on the realized next-period productivity. The labor boundary signs from part C apply with $f_l$ replaced by $a_t f_l$. For a well-behaved optimum the capital path must also satisfy the appropriate expected transversality condition, typically
\[
\lim_{T\to\infty}\mathbb{E}_0\left[\beta^T u_c(c_T,l_T)k_{T+1}\right]=0.
\]
As in the deterministic case, first-order conditions alone do not establish global optimality under the stated curvature assumptions.

\href{ps1.tex}{Back to the exercises}
